% ================================================================
% Consolidated Mathematics Notes — Main Document
% ================================================================
% Structure:
%   main.tex                 <- you are here
%   parts/
%     p0_psle/               <- PSLE
%     p1_emath/              <- O-Level E Math
%     p2_amath/              <- O-Level A Math
%     p3_h2math/             <- A-Level H2 Math
%     p4_olympiad/           <- Math Olympiad
%     p5_discrete/           <- CS1231S Discrete Structures
%     p6_calculus/           <- MA1521 + MA2104
%     p7_linalg/             <- MA2001 + MA2101
%     p8_analysis/           <- MA2108
%     p9_prob/               <- ST2334
%     p10_stochastic/        <- MA3238 + MA4251
%     p11_rl/                <- MA4275
% ================================================================

\documentclass[11pt,a4paper]{article}

% --- Core packages ---
\usepackage[utf8]{inputenc}
\usepackage[T1]{fontenc}
\usepackage{lmodern}
\usepackage[margin=0.9in]{geometry}
\usepackage{amsmath, amssymb, amsthm}
\usepackage{mathtools}
\usepackage{bm}
\usepackage{hyperref}
\usepackage{enumitem}
\usepackage{fancyhdr}
\usepackage{titlesec}
\usepackage{tocloft}
\usepackage{parskip}
\usepackage{booktabs}
\usepackage{array}
\usepackage{pgfplots}
\usepackage{tikz}
\usepackage{xcolor}
\usepackage{mdframed}
\usepackage{thmtools}
\pgfplotsset{compat=1.18}
\usetikzlibrary{positioning, arrows.meta, calc, patterns}
\setlength{\headheight}{13.6pt}
\addtolength{\topmargin}{-1.6pt}

% --- Theorem environments ---
\declaretheoremstyle[
  headfont=\bfseries,
  bodyfont=\normalfont,
  headpunct={.},
  postheadspace=0.5em
]{thmstyle}

\declaretheorem[style=thmstyle, name=Theorem,     numberwithin=section]{theorem}
\declaretheorem[style=thmstyle, name=Lemma,       sibling=theorem]{lemma}
\declaretheorem[style=thmstyle, name=Corollary,   sibling=theorem]{corollary}
\declaretheorem[style=thmstyle, name=Proposition, sibling=theorem]{proposition}
\declaretheorem[style=thmstyle, name=Definition,  sibling=theorem]{definition}
\declaretheorem[style=thmstyle, name=Example,     sibling=theorem]{example}
\declaretheorem[style=thmstyle, name=Remark,      sibling=theorem]{remark}

% --- Coloured theorem boxes ---
\definecolor{thmblue}{RGB}{220,232,248}
\definecolor{defgreen}{RGB}{225,243,228}
\definecolor{exorange}{RGB}{253,244,220}
\surroundwithmdframed[backgroundcolor=thmblue,  linewidth=0pt, innertopmargin=4pt, innerbottommargin=4pt]{theorem}
\surroundwithmdframed[backgroundcolor=defgreen, linewidth=0pt, innertopmargin=4pt, innerbottommargin=4pt]{definition}
\surroundwithmdframed[backgroundcolor=exorange, linewidth=0pt, innertopmargin=4pt, innerbottommargin=4pt]{example}

% --- Common math shorthands ---
\newcommand{\R}{\mathbb{R}}
\newcommand{\N}{\mathbb{N}}
\newcommand{\Z}{\mathbb{Z}}
\newcommand{\Q}{\mathbb{Q}}
\newcommand{\C}{\mathbb{C}}
\newcommand{\F}{\mathbb{F}}
\newcommand{\E}{\mathbb{E}}
\newcommand{\Prob}{\mathbb{P}}
\newcommand{\Var}{\operatorname{Var}}
\newcommand{\Cov}{\operatorname{Cov}}
\newcommand{\rank}{\operatorname{rank}}
\newcommand{\nullity}{\operatorname{nullity}}
\newcommand{\tr}{\operatorname{tr}}
\newcommand{\norm}[1]{\left\lVert #1 \right\rVert}
\newcommand{\abs}[1]{\left\lvert #1 \right\rvert}
\newcommand{\inner}[2]{\langle #1,\, #2 \rangle}
\newcommand{\eps}{\varepsilon}
\newcommand{\dd}{\,\mathrm{d}}
\newcommand{\pd}[2]{\frac{\partial #1}{\partial #2}}

% --- Page style ---
\pagestyle{fancy}
\fancyhf{}
\fancyhead[L]{\leftmark}
\fancyhead[R]{Mathematics Notes}
\fancyfoot[C]{\thepage}
\renewcommand{\headrulewidth}{0.4pt}

% --- Section formatting ---
\titleformat{\section}[block]{\Large\bfseries}{}{0em}{}
\newcommand{\sectionbreak}{\newpage}
\titleformat{\subsection}[block]{\large\bfseries}{}{0em}{}
\titleformat{\subsubsection}[block]{\normalsize\bfseries}{}{0em}{}

% --- Hyperlinks ---
\hypersetup{
  colorlinks=true,
  linkcolor=blue!60!black,
  urlcolor=blue!60!black,
  citecolor=blue!60!black,
}

% --- Title ---
\title{\textbf{Consolidated Mathematics Notes}\\[0.5em]
  \Large PSLE $\to$ O-Level $\to$ A-Level $\to$ Olympiad $\to$ University}
\author{Jeriel Chan}
\date{\today}

\begin{document}

\maketitle
\tableofcontents
\newpage

% ================================================================
%  PART 0 — Arithmetical Foundations (PSLE)
% ================================================================
\part{Arithmetical Foundations}
% ==============================================================
% Numbers and the Number System
% ==============================================================

\section{Numbers and the Number System}

\subsection{Natural, integer, rational, real numbers}

% TODO

\subsection{Primes, factors, HCF/LCM}

% TODO

\subsection{Standard form and estimation}

% TODO

% ==============================================================
% Fractions, Decimals, Percentages
% ==============================================================

\section{Fractions, Decimals, Percentages}

\subsection{Equivalent fractions and operations}

% TODO

\subsection{Percentage change, profit and loss}

% TODO

\subsection{Reverse percentage problems}

% TODO

% ==============================================================
% Ratio, Rate, and Proportion
% ==============================================================

\section{Ratio, Rate, and Proportion}

\subsection{Part-whole and part-part ratios}

% TODO

\subsection{Speed, distance, time}

% TODO

\subsection{Direct and inverse proportion}

% TODO

% ==============================================================
% Basic Geometry and Mensuration
% ==============================================================

\section{Basic Geometry and Mensuration}

\subsection{Angles: complementary, supplementary, vertically opposite}

% TODO

\subsection{Area and perimeter of 2D shapes}

% TODO

\subsection{Volume and surface area of 3D solids}

% TODO

% ==============================================================
% Data Handling
% ==============================================================

\section{Data Handling}

\subsection{Mean, median, mode}

% TODO

\subsection{Tables, bar charts, line graphs, pie charts}

% TODO

\subsection{Interpretation and problem solving}

% TODO


% ================================================================
%  PART I — O-Level Elementary Mathematics
% ================================================================
\part{O-Level Elementary Mathematics}
% ==============================================================
% Algebraic Manipulation
% ==============================================================

\section{Algebraic Manipulation}

\subsection{Expansion and factorisation}

% TODO

\subsection{Quadratic expressions and equations}

% TODO

\subsection{Algebraic fractions and surds}

% TODO

% ==============================================================
% Functions and Graphs
% ==============================================================

\section{Functions and Graphs}

\subsection{Linear, quadratic, cubic, reciprocal}

% TODO

\subsection{Exponential and modulus graphs}

% TODO

\subsection{Sketching and transformation}

% TODO

% ==============================================================
% Coordinate Geometry
% ==============================================================

\section{Coordinate Geometry}

\subsection{Gradient and intercept form}

% TODO

\subsection{Distance, midpoint, perpendicular bisector}

% TODO

\subsection{Collinearity}

% TODO

% ==============================================================
% Trigonometry
% ==============================================================

\section{Trigonometry}

\subsection{Right-angled triangles, SOH-CAH-TOA}

% TODO

\subsection{Sine and cosine rule}

% TODO

\subsection{Angles of elevation and depression}

% TODO

% ==============================================================
% Statistics and Probability
% ==============================================================

\section{Statistics and Probability}

\subsection{Frequency tables, cumulative frequency, box plots}

% TODO

\subsection{Mean, variance from grouped data}

% TODO

\subsection{Classical probability, tree diagrams, Venn diagrams}

% TODO

% ==============================================================
% Mensuration
% ==============================================================

\section{Mensuration}

\subsection{Arc length, sector area, segment}

% TODO

\subsection{Pyramid, cone, sphere volumes}

% TODO

\subsection{Composite solids}

% TODO


% ================================================================
%  PART II — O-Level Additional Mathematics
% ================================================================
\part{O-Level Additional Mathematics}
% ==============================================================
% Polynomials
% ==============================================================

\section{Polynomials}

\subsection{Remainder and factor theorems}

% TODO

\subsection{Cubic factorisation}

% TODO

\subsection{Partial fractions}

% TODO

% ==============================================================
% Exponential and Logarithmic Functions
% ==============================================================

\section{Exponential and Logarithmic Functions}

\subsection{Laws of logarithms}

% TODO

\subsection{Equations and inequalities}

% TODO

\subsection{Graphs and modelling}

% TODO

% ==============================================================
% Trigonometry (Advanced)
% ==============================================================

\section{Trigonometry (Advanced)}

\subsection{Compound and double angle formulae}

% TODO

\subsection{R-formula}

% TODO

\subsection{General solutions and graphs}

% TODO

% ==============================================================
% Coordinate Geometry (Advanced)
% ==============================================================

\section{Coordinate Geometry (Advanced)}

\subsection{Circles: general and standard form}

% TODO

\subsection{Points of intersection, tangent and normal}

% TODO

\subsection{Introduction to conics}

% TODO

% ==============================================================
% Binomial Theorem
% ==============================================================

\section{Binomial Theorem}

\subsection{Binomial expansion for positive integers}

% TODO

\subsection{General term and coefficient finding}

% TODO

\subsection{Applications and approximations}

% TODO

% ==============================================================
% Introductory Calculus
% ==============================================================

\section{Introductory Calculus}

\subsection{Differentiation: product, quotient, chain rules}

% TODO

\subsection{Integration as anti-differentiation}

% TODO

\subsection{Definite integrals and area under curves}

% TODO

% ==============================================================
% Sequences and Series
% ==============================================================

\section{Sequences and Series}

\subsection{Arithmetic and geometric progressions}

% TODO

\subsection{Sum to infinity of GP}

% TODO

\subsection{Sigma notation}

% TODO


% ================================================================
%  PART III — A-Level H2 Mathematics
% ================================================================
\part{A-Level H2 Mathematics}
% ==============================================================
% Functions
% ==============================================================

\section{Functions}

\subsection{Domain, range, inverse, composite}

% TODO

\subsection{Odd and even functions}

% TODO

\subsection{Modulus functions and equations}

% TODO

% ==============================================================
% Sequences and Series (H2)
% ==============================================================

\section{Sequences and Series (H2)}

\subsection{Method of differences}

% TODO

\subsection{Convergence, limit of sum}

% TODO

\subsection{Maclaurin and Taylor series (intro)}

% TODO

% ==============================================================
% Vectors
% ==============================================================

\section{Vectors}

\subsection{2D and 3D vectors, dot and cross product}

% TODO

\subsection{Lines and planes in 3D}

% TODO

\subsection{Distances, intersections, angles}

% TODO

% ==============================================================
% Complex Numbers
% ==============================================================

\section{Complex Numbers}

\subsection{Cartesian, polar, and exponential forms}

% TODO

\subsection{Loci in the Argand diagram}

% TODO

\subsection{Roots of unity and de Moivre's theorem}

% TODO

% ==============================================================
% Calculus (H2)
% ==============================================================

\section{Calculus (H2)}

\subsection{Advanced integration techniques}

% TODO

\subsection{Volumes of revolution}

% TODO

\subsection{Differential equations: separable, integrating factor}

% TODO

% ==============================================================
% Permutations and Combinations
% ==============================================================

\section{Permutations and Combinations}

\subsection{Fundamental counting principle}

% TODO

\subsection{Arrangements with restrictions}

% TODO

\subsection{Combinations and selections}

% TODO

% ==============================================================
% Probability (H2)
% ==============================================================

\section{Probability (H2)}

\subsection{Conditional probability, Bayes' theorem}

% TODO

\subsection{Binomial, Poisson, Normal distributions}

% TODO

\subsection{Continuity correction}

% TODO

% ==============================================================
% Hypothesis Testing
% ==============================================================

\section{Hypothesis Testing}

\subsection{z-tests and t-tests}

% TODO

\subsection{p-values and critical regions}

% TODO

\subsection{Correlation and linear regression}

\subsubsection*{Ordinary Least Squares (OLS) Regression}

\begin{definition}
In \textbf{multiple linear regression}, we model a response variable $y$ as a linear function of $p$ predictors:
\[
\mathbf{y} = \mathbf{X}\boldsymbol{\beta} + \boldsymbol{\varepsilon}, \qquad \boldsymbol{\varepsilon} \sim (0, \sigma^2 \mathbf{I})
\]
where $\mathbf{y} \in \mathbb{R}^n$, $\mathbf{X} \in \mathbb{R}^{n \times p}$ is the \emph{design matrix}, and $\boldsymbol{\beta} \in \mathbb{R}^p$ is the coefficient vector. The \textbf{Ordinary Least Squares} (OLS) estimator minimises the sum of squared residuals:
\[
\hat{\boldsymbol{\beta}}_{\mathrm{OLS}} = \arg\min_{\boldsymbol{\beta}} \|\mathbf{y} - \mathbf{X}\boldsymbol{\beta}\|^2 = (\mathbf{X}^\top \mathbf{X})^{-1}\mathbf{X}^\top \mathbf{y}
\]
\end{definition}

\begin{remark}
The closed-form solution $(\mathbf{X}^\top \mathbf{X})^{-1}$ exists \emph{only if} $\mathbf{X}^\top\mathbf{X}$ is invertible, which requires $\mathbf{X}$ to have full column rank. This is precisely the condition violated by perfect multicollinearity.
\end{remark}

\subsubsection*{Derivation of the OLS Estimator}

The OLS estimator is found by minimising the \textbf{sum of squared residuals} (the loss function):
\[
\mathcal{L}(\boldsymbol{\beta}) = (\mathbf{y} - \mathbf{X}\boldsymbol{\beta})^\top(\mathbf{y} - \mathbf{X}\boldsymbol{\beta})
\]

\textbf{Step 1 — Expand the loss.} Using $(a-b)^\top(a-b) = a^\top a - 2b^\top a + b^\top b$:
\[
\mathcal{L}(\boldsymbol{\beta}) = \mathbf{y}^\top\mathbf{y} - 2\boldsymbol{\beta}^\top\mathbf{X}^\top\mathbf{y} + \boldsymbol{\beta}^\top\mathbf{X}^\top\mathbf{X}\boldsymbol{\beta}
\]

\textbf{Step 2 — Differentiate and set to zero.} Taking the gradient with respect to $\boldsymbol{\beta}$:
\[
\frac{\partial \mathcal{L}}{\partial \boldsymbol{\beta}} = -2\mathbf{X}^\top\mathbf{y} + 2\mathbf{X}^\top\mathbf{X}\boldsymbol{\beta} = \mathbf{0}
\]

\textbf{Step 3 — Normal equations.} Rearranging gives:
\[
\mathbf{X}^\top\mathbf{X}\,\boldsymbol{\beta} = \mathbf{X}^\top\mathbf{y}
\]

\textbf{Step 4 — Solve.} Provided $\mathbf{X}^\top\mathbf{X}$ is invertible, multiply both sides on the \emph{left} by $(\mathbf{X}^\top\mathbf{X})^{-1}$:
\[
\hat{\boldsymbol{\beta}}_{\mathrm{OLS}} = (\mathbf{X}^\top\mathbf{X})^{-1}\mathbf{X}^\top\mathbf{y}
\]

\begin{remark}
Note on order: writing $\mathbf{X}^\top\mathbf{y}\,(\mathbf{X}^\top\mathbf{X})^{-1}$ is \textbf{wrong} — matrix multiplication is not commutative. The inverse must appear on the left.
\end{remark}

\subsubsection*{The No Perfect Multicollinearity Assumption}

\begin{definition}
The \textbf{no perfect multicollinearity} assumption requires that no predictor variable is an exact linear combination of the others. Equivalently, the design matrix $\mathbf{X}$ must have full column rank: $\mathrm{rank}(\mathbf{X}) = p$.
\end{definition}

\begin{remark}
If perfect multicollinearity holds (e.g.\ $X_3 = 2X_1 + X_2$ exactly), then $\mathbf{X}^\top\mathbf{X}$ is singular and its inverse does not exist. OLS estimation breaks down completely --- there are infinitely many coefficient vectors that fit the data equally well, so no unique solution can be identified.
\end{remark}

\begin{example}
Suppose a model predicts income using both $X_1 = \text{age in months}$ and $X_2 = \text{age in years}$. Since $X_1 = 12 X_2$ exactly, the design matrix is rank-deficient. The model cannot separately estimate the effect of one month versus one year --- they always move together. \textbf{Fix:} drop one of the two variables.
\end{example}

\subsubsection*{Common Causes}

\begin{remark}
The three most frequent sources of perfect multicollinearity are:
\begin{itemize}
    \item \textbf{Dummy variable trap:} Including all $k$ categories of a categorical variable as dummy variables alongside an intercept. Since the dummies sum to 1 (the intercept column), perfect collinearity is introduced. \textbf{Fix:} drop one reference category.
    \item \textbf{Duplicate variables:} Including the same measurement twice in different units (e.g.\ area in $\text{m}^2$ and area in $\text{cm}^2$).
    \item \textbf{Constant predictors:} A predictor with zero variance (same value for all observations) is perfectly collinear with the intercept.
\end{itemize}
\end{remark}

\subsubsection*{Perfect vs.\ Imperfect Multicollinearity}

\begin{definition}
\textbf{Imperfect (near) multicollinearity} occurs when predictors are highly correlated but not in an exact linear relationship. OLS still produces a unique solution, but coefficient estimates become highly sensitive to small changes in the data and standard errors inflate, making inference unreliable.
\end{definition}

\begin{remark}
The \textbf{Variance Inflation Factor} (VIF) measures the severity of imperfect multicollinearity for predictor $j$:
\[
\mathrm{VIF}_j = \frac{1}{1 - R_j^2}
\]
where $R_j^2$ is the $R^2$ from regressing $X_j$ on all other predictors. A rule of thumb: $\mathrm{VIF}_j > 10$ signals problematic multicollinearity.
\end{remark}

\subsubsection*{Heteroscedasticity}

\begin{definition}
The OLS error term $\varepsilon_i$ is \textbf{homoscedastic} if its conditional variance is constant across all observations:
\[
\mathrm{Var}(\varepsilon_i \mid X_i) = \sigma^2 \quad \forall\, i
\]
It is \textbf{heteroscedastic} if this variance instead depends on $X_i$:
\[
\mathrm{Var}(\varepsilon_i \mid X_i) = \sigma_i^2 \quad \text{(varies with } i\text{)}
\]
\end{definition}

\textbf{Variance structure.} Common parametric forms for $\sigma_i^2$ include:
\[
\sigma_i^2 = \sigma^2 \cdot X_i, \qquad \sigma_i^2 = \sigma^2 \cdot X_i^2, \qquad \sigma_i^2 = \exp(\alpha + \beta X_i)
\]
The subscript $i$ is the essential signal: each observation carries its own variance rather than a single pooled value.

\textbf{Visual signature.} A residual-vs-fitted plot showing a \emph{funnel} (spread growing with fitted values) is the classic indicator of heteroscedasticity. A reverse funnel is equally diagnostic.

\begin{remark}[Consequences for OLS]
Under heteroscedasticity:
\begin{itemize}
    \item \textbf{Coefficients remain unbiased} — $\hat{\boldsymbol{\beta}}_{\mathrm{OLS}}$ still converges to the true $\boldsymbol{\beta}$ in expectation.
    \item \textbf{Gauss--Markov fails} — OLS is no longer the Best Linear Unbiased Estimator (BLUE); more efficient estimators (e.g.\ WLS) exist.
    \item \textbf{Standard errors are biased} — standard errors computed under homoscedasticity are typically \emph{underestimated}, inflating $t$-statistics.
    \item \textbf{Inference breaks down} — because standard errors feed into $t$-tests, $F$-tests, and confidence intervals, all classical hypothesis tests become unreliable and may produce spurious rejections of $H_0$.
\end{itemize}
\end{remark}

\textbf{Detection tests.}
\begin{itemize}
    \item \textbf{Visual:} Plot residuals vs.\ fitted values; look for funnel or fan shapes.
    \item \textbf{Breusch--Pagan test:} Regress squared residuals $\hat{\varepsilon}_i^2$ on the regressors. The LM test statistic follows $\chi^2(k)$ under $H_0$ (homoscedasticity).
    \item \textbf{White test:} Like Breusch--Pagan but also includes squares and cross-products of regressors, making it more robust to functional-form misspecification.
    \item \textbf{Goldfeld--Quandt test:} Split the sample into two groups; compare their residual variances via an $F$-test.
\end{itemize}

\begin{remark}[Remedies]
\begin{itemize}
    \item \textbf{Heteroscedasticity-consistent (HC) standard errors} (robust SEs): correct standard error estimates without assuming constant variance. Common variants are HC0 (White 1980), HC1 (degrees-of-freedom corrected), and HC3 (recommended for small samples).
    \item \textbf{Weighted Least Squares (WLS):} down-weights high-variance observations. Efficient when the variance structure $\sigma_i^2$ is known or estimable.
    \item \textbf{Generalised Least Squares (GLS):} full generalisation of WLS for an arbitrary variance--covariance structure.
    \item \textbf{Variance-stabilising transformations:} apply $\ln(Y)$ or a power transformation to the dependent variable to stabilise the spread of residuals.
\end{itemize}
\end{remark}

\subsubsection*{Joint Hypothesis Testing: the $F$-Statistic in Regression}

A \textbf{$t$-test} checks one coefficient at a time. An \textbf{$F$-test} checks several coefficients simultaneously. To test $H_0\colon \beta_1 = 0$ \emph{and} $\beta_2 = 0$ together, combine both $t$-statistics into a single score.

\begin{theorem}[Joint $F$-statistic, two restrictions]
Let $t_1 = \hat\beta_1 / \widehat{\mathrm{SE}}_1$ and $t_2 = \hat\beta_2 / \widehat{\mathrm{SE}}_2$ be the individual $t$-statistics, and let $\hat\rho \equiv \hat\rho_{t_1,t_2}$ be the estimated correlation between them. Then the joint $F$-statistic is
\[
F = \frac{1}{2} \cdot \frac{t_1^2 + t_2^2 - 2\hat\rho\, t_1 t_2}{1 - \hat\rho^2}.
\]
Under $H_0$, in large samples, $F \sim F(2,\infty)$.
\end{theorem}

\textbf{Anatomy of the formula.}
\begin{itemize}
    \item $t_1^2 + t_2^2$: raw evidence. Squaring ensures large negative and large positive $t$-values both count.
    \item $-2\hat\rho\, t_1 t_2$: removes double-counting. If $t_1$ and $t_2$ move together ($\hat\rho > 0$), they share information; this term subtracts the overlap.
    \item $1 - \hat\rho^2$: rescaling. Even after removing the overlap in the numerator, correlated variables do not spread like independent ones — this denominator adjusts for the ``compressed'' joint geometry.
    \item $\frac{1}{2}$: divides by the number of restrictions (2), making the result comparable to $F(2,\cdot)$.
\end{itemize}

\begin{remark}[Special case: uncorrelated $t$-statistics]
When $\hat\rho = 0$ (regressors uncorrelated), the formula collapses to
\[
F = \tfrac{1}{2}(t_1^2 + t_2^2).
\]
This is simply the average of the squared $t$-statistics. In large samples, $t_1, t_2 \approx N(0,1)$ independently under $H_0$, so $t_1^2 + t_2^2 \sim \chi^2(2)$ and $F \sim F(2,\infty)$.
\end{remark}

\begin{remark}[General: $q$ restrictions]
Testing $q$ restrictions simultaneously gives $F \sim F(q,\infty)$ in large samples. The $q$ in the numerator degrees of freedom is precisely the number of restrictions being tested.
\end{remark}

\subsubsection*{Mahalanobis Distance and the $F$-Formula}

The joint $F$-statistic is a scaled \textbf{Mahalanobis distance} — the correct notion of distance when two variables are correlated.

Under $H_0$, $(t_1, t_2)^\top$ has covariance matrix
\[
\Sigma = \begin{pmatrix} 1 & \rho \\ \rho & 1 \end{pmatrix},
\]
where each $t$-stat has variance 1 (both are approximately $N(0,1)$) and they share correlation $\rho$ because $\hat\beta_1, \hat\beta_2$ come from the same regression data.

Inverting a $2\times 2$ matrix $\begin{pmatrix}a&b\\c&d\end{pmatrix}$ via $\frac{1}{ad-bc}\begin{pmatrix}d&-b\\-c&a\end{pmatrix}$, with $\det(\Sigma) = 1 - \rho^2$:
\[
\Sigma^{-1} = \frac{1}{1-\rho^2}\begin{pmatrix} 1 & -\rho \\ -\rho & 1 \end{pmatrix}.
\]

The Mahalanobis squared distance from the origin is $\mathbf{t}^\top \Sigma^{-1} \mathbf{t}$:
\[
\mathbf{t}^\top \Sigma^{-1} \mathbf{t}
= \frac{1}{1-\rho^2}(t_1, t_2)\begin{pmatrix}1 & -\rho \\ -\rho & 1\end{pmatrix}\begin{pmatrix}t_1\\t_2\end{pmatrix}
= \frac{t_1^2 + t_2^2 - 2\rho t_1 t_2}{1-\rho^2}.
\]
The $-2\rho t_1 t_2$ arises because the off-diagonal $-\rho$ entries appear in both rows of the matrix (symmetry), contributing $-\rho t_1 t_2$ twice. The $1-\rho^2$ is the determinant of $\Sigma$, falling out of the matrix inversion.

Dividing by 2 (the number of restrictions) gives the $F$-statistic $F = \frac{1}{2}\mathbf{t}^\top \hat\Sigma^{-1}\mathbf{t}$.

\begin{remark}[Geometric intuition: circles vs.\ ellipses]
When $\rho = 0$, the equal-distance contours of $(t_1,t_2)$ from the origin are circles. When $\rho \neq 0$, the cloud of possible values stretches along the diagonal — equal-distance contours become ellipses. Points along $t_1 \approx t_2$ are not genuinely independent evidence; they lie in a correlated direction. The Mahalanobis formula measures distance in this elliptical geometry, not the naive circular one. As $|\rho| \to 1$, the ellipse degenerates and the test becomes unstable ($1-\rho^2 \to 0$).
\end{remark}

\begin{remark}[Why $t$-statistics can be correlated in regression]
The $t$-stats $t_j = \hat\beta_j / \widehat{\mathrm{SE}}_j$ share the same sample data. If regressors $X_1$ and $X_2$ are correlated with each other, the coefficient estimates $\hat\beta_1$ and $\hat\beta_2$ are correlated too — they trade off against each other to fit the same outcomes. This is why $\hat\rho \neq 0$ is the norm, not the exception.
\end{remark}


% ================================================================
%  PART IV — Mathematical Olympiad Techniques
% ================================================================
\part{Mathematical Olympiad Techniques}
% ==============================================================
% Number Theory (Olympiad)
% ==============================================================

\section{Number Theory (Olympiad)}

\subsection{Divisibility and modular arithmetic}

% TODO

\subsection{Euler's theorem and CRT}

% TODO

\subsection{Legendre symbol, quadratic reciprocity}

% TODO

% ==============================================================
% Combinatorics (Olympiad)
% ==============================================================

\section{Combinatorics (Olympiad)}

\subsection{Principle of inclusion-exclusion}

\begin{theorem}[Inclusion-Exclusion Principle]
For finite sets $A_1, A_2, \ldots, A_n$:
\[
\left|A_1 \cup A_2 \cup \cdots \cup A_n\right|
= \sum_i |A_i| - \sum_{i<j}|A_i \cap A_j| + \sum_{i<j<k}|A_i \cap A_j \cap A_k| - \cdots
\]
For two sets this simplifies to $|A \cup B| = |A| + |B| - |A \cap B|$.
\end{theorem}

\subsubsection*{Summing Multiples with Inclusion-Exclusion}

A standard application is summing all multiples of $a$ or $b$ up to some limit.

\begin{example}[Sum of multiples of 3 or 5 below 1000]
Find the sum of all natural numbers below 1000 that are multiples of 3 or 5.

\textbf{Multiples of 3 below 1000}: these are $3, 6, 9, \ldots, 999$. There are $\lfloor 999/3 \rfloor = 333$ terms, so their sum is
\[
3(1+2+\cdots+333) = 3 \cdot \frac{333 \cdot 334}{2} = 166833.
\]

\textbf{Multiples of 5 below 1000}: $5, 10, \ldots, 995$. There are $\lfloor 999/5 \rfloor = 199$ terms, so their sum is
\[
5(1+2+\cdots+199) = 5 \cdot \frac{199 \cdot 200}{2} = 99500.
\]

\textbf{Multiples of both} (i.e.\ multiples of $\text{lcm}(3,5)=15$): $15, 30, \ldots, 990$. There are $\lfloor 999/15 \rfloor = 66$ terms, so their sum is
\[
15(1+2+\cdots+66) = 15 \cdot \frac{66 \cdot 67}{2} = 33165.
\]

By inclusion-exclusion:
\[
166833 + 99500 - 33165 = \mathbf{233168}.
\]
\end{example}

\begin{remark}[General template]
To sum all multiples of $d$ below $N$: let $k = \lfloor (N-1)/d \rfloor$. The sum is $d \cdot k(k+1)/2$. For multiples of $a$ or $b$ below $N$, subtract the multiples of $\text{lcm}(a,b)$.
\end{remark}

\subsection{Generating functions}

% TODO

\subsection{Bijections and double counting}

% TODO

\subsection{Combinatorial game theory}

\subsubsection*{Two-Pile Threshold Games}

Consider the following two-player game: two piles of stones of sizes $a$ and $b$ ($a \le b$). A move consists of removing at least one stone from each pile, with the \emph{total} number removed equal to $\min(a,b)$. The player who cannot move loses.

\textbf{Key reduction.} From a position $(a,b)$ with $a \le b$, exactly $a$ stones are removed in total (since $a = \min(a,b)$). The total remaining is $(a+b)-a = b$. So every move from $(a,b)$ lands on a position whose pile sum is exactly $b$. Equivalently, the reachable positions from $(a,b)$ are exactly:
\[
\{(c,\, b-c) : 1 \le c \le \min(a, \lfloor b/2 \rfloor)\}.
\]
For fixed $b$, increasing $a$ only \emph{adds} reachable moves — it never removes them. This gives a monotonicity principle:

\begin{theorem}[Column threshold]
For fixed $b$, there exists a threshold $r(b)$ such that $(a,b)$ is a losing position if and only if $a < r(b)$.
\end{theorem}

\subsubsection*{The Lowbit Threshold}

\begin{theorem}
The threshold for column $b$ is
\[
r(b) = \operatorname{lowbit}(b+1) = 2^{v_2(b+1)},
\]
i.e.\ the largest power of $2$ dividing $b+1$. The position $(a,b)$ (with $a \le b$) is \textbf{losing} if and only if
\[
a < \operatorname{lowbit}(b+1).
\]
\end{theorem}

\begin{proof}[Proof sketch]
Let $r = \operatorname{lowbit}(b+1)$. A child $(c, b-c)$ is losing iff $c < \operatorname{lowbit}(b-c+1)$.

\emph{Case $1 \le c < r$}: Then $\operatorname{lowbit}((b+1)-c) = \operatorname{lowbit}(c) \le c$, so the child is \emph{not} losing (it is a winning position).

\emph{Case $c = r$}: Since $r = \operatorname{lowbit}(b+1)$, subtracting $r$ from $b+1$ eliminates the lowest set bit, making $\operatorname{lowbit}(b+1-r) > r$. So $r < \operatorname{lowbit}(b-r+1)$, meaning $(r, b-r)$ \emph{is} losing.

Therefore: the first losing child is at $c = r$. Any position with $a < r$ cannot reach it, so $(a,b)$ is losing. Any position with $a \ge r$ can move to $(r, b-r)$ (losing for the opponent), so it is winning.
\end{proof}

\begin{example}[Small cases]
\begin{center}
\begin{tabular}{ccll}
$b$ & $b+1$ & $\operatorname{lowbit}(b+1)$ & Losing positions $(a,b)$, $a \le b$ \\
\hline
3 & 4 & 4 & all of $(1,3),(2,3),(3,3)$ \\
5 & 6 & 2 & only $(1,5)$ \\
7 & 8 & 8 & all of $(1,7),\ldots,(7,7)$ \\
11 & 12 & 4 & $(1,11),(2,11),(3,11)$ \\
\end{tabular}
\end{center}
The ``all-losing columns'' occur when $\operatorname{lowbit}(b+1) > b$, i.e.\ when $b+1$ is itself a power of 2, so $b = 2^k - 1$: columns $1, 3, 7, 15, 31, \ldots$
\end{example}

\subsubsection*{Counting Losing Positions}

\begin{theorem}[Number of ordered losing pairs $L(n)$]
For ordered pairs $(a,b)$ with $1 \le a,b \le n$:
\[
L(n) = 2\bigl(T(n+1) - 1\bigr) - 2n - \lfloor \log_2(n+1) \rfloor,
\]
where $T(N) = \sum_{m=1}^{N} \operatorname{lowbit}(m)$.
\end{theorem}

\begin{proof}[Derivation]
For each column $b$, the number of losing values $a \in \{1,\ldots,b\}$ is $\operatorname{lowbit}(b+1) - 1$. For ordered pairs, off-diagonal positions $(a,b)$ and $(b,a)$ are counted separately; the diagonal $(b,b)$ is losing exactly when $b < \operatorname{lowbit}(b+1)$, i.e.\ when $b = 2^k - 1$. Summing and adjusting for symmetry gives the formula above.
\end{proof}

\begin{remark}[Fast computation]
The sum $T(N)$ satisfies the recurrence $T(2k) = k + 2T(k)$ and $T(2k+1) = T(2k)+1$ (see the 2-adic valuation subsection in the number theory chapter), enabling $O(\log N)$ evaluation. This makes $L(n)$ computable in logarithmic time for huge $n$.
\end{remark}

\begin{example}
$L(7) = 21$, $L(49) = 221$, $L(7^{17}) = 10784223938983273$.
\end{example}

% ==============================================================
% Geometry (Olympiad)
% ==============================================================

\section{Geometry (Olympiad)}

\subsection{Circle theorems and power of a point}

% TODO

\subsection{Trigonometric methods}

% TODO

\subsection{Projective geometry and inversions}

% TODO

% ==============================================================
% Algebra (Olympiad)
% ==============================================================

\section{Algebra (Olympiad)}

\subsection{AM-GM, Cauchy-Schwarz, rearrangement}

% TODO

\subsection{Functional equations}

% TODO

\subsection{Polynomial tricks and symmetric functions}

% TODO

% ==============================================================
% Proof Strategies
% ==============================================================

\section{Proof Strategies}

\subsection{Invariants and monovariants}

% TODO

\subsection{Extremal principle}

% TODO

\subsection{Pigeonhole principle}

% TODO


% ================================================================
%  PART V — Discrete Structures  (CS1231S)
% ================================================================
\part{Discrete Structures}
% ==============================================================
% Logic and Proofs
% ==============================================================

\section{Logic and Proofs}

\subsection{Propositional and predicate logic}

\subsubsection*{Logical Connectives and XOR}

\begin{definition}
The standard propositional connectives are $\neg$ (NOT), $\land$ (AND), $\lor$ (OR), $\Rightarrow$ (implication), and $\Leftrightarrow$ (biconditional). The \textbf{exclusive or} (XOR) of propositions $p$ and $q$, written $p \oplus q$, is true exactly when $p$ and $q$ have different truth values:
\[
p \oplus q \;=\; (p \lor q) \land \neg(p \land q).
\]
\end{definition}

\begin{remark}[XOR truth table]
\begin{center}
\begin{tabular}{c|c|c}
$p$ & $q$ & $p \oplus q$ \\
\hline
0 & 0 & 0 \\
0 & 1 & 1 \\
1 & 0 & 1 \\
1 & 1 & 0 \\
\end{tabular}
\end{center}
XOR outputs 1 exactly when its inputs differ. Equivalently, $p \oplus q \equiv p + q \pmod{2}$: it is binary addition without carry.
\end{remark}

\begin{example}[Bitwise XOR]
XOR extends bitwise to binary strings of equal length, applying the operation bit by bit:
\[
101 \oplus 110 = 011
\]
Check: $1\oplus1=0$, $0\oplus1=1$, $1\oplus0=1$.

In quantum computing, expressions like $y \leftarrow x \oplus s$ mean ``flip the bits of $x$ wherever $s$ has a 1.'' For instance, $101 \oplus 110 = 011$.
\end{example}

\begin{remark}[Key properties of $\oplus$]
\begin{itemize}
    \item \textbf{Self-inverse:} $x \oplus x = 0$ for any bit string $x$.
    \item \textbf{Identity:} $x \oplus 0 = x$.
    \item \textbf{Commutativity:} $x \oplus y = y \oplus x$.
    \item \textbf{Associativity:} $(x \oplus y) \oplus z = x \oplus (y \oplus z)$.
\end{itemize}
Self-inverse is especially useful: if $y = x \oplus s$, then $y \oplus s = x$ (XOR with the same value recovers the original).
\end{remark}

\subsection{Direct, contrapositive, contradiction}

% TODO

\subsection{Mathematical induction (weak and strong)}

% TODO

% ==============================================================
% Set Theory and Functions
% ==============================================================

\section{Set Theory and Functions}

\subsection{Set operations, cardinality, power sets}

% TODO

\subsection{Injective, surjective, bijective}

% TODO

\subsection{Cantor's theorem, countability}

% TODO

% ==============================================================
% Relations
% ==============================================================

\section{Relations}

\subsection{Equivalence relations and partitions}

% TODO

\subsection{Partial orders, Hasse diagrams}

% TODO

\subsection{Well-ordering principle}

% TODO

% ==============================================================
% Combinatorics (Discrete)
% ==============================================================

\section{Combinatorics (Discrete)}

\subsection{Counting: permutations and combinations}

% TODO

\subsection{Binomial theorem and Pascal's identity}

% TODO

\subsection{Recurrence relations and closed forms}

% TODO

\subsection{Digit DP}

\textbf{Digit DP} is a technique for counting integers in a range $[1, N]$ that satisfy some property that can be decomposed digit by digit. The key idea is to maintain a DP state that tracks:
\begin{itemize}
    \item how many digits are left to place,
    \item any running modular conditions (e.g.\ divisibility by $k$),
    \item any parity or membership constraints on digit counts,
    \item whether the prefix built so far is still ``tight'' (equal to $N$'s prefix).
\end{itemize}

\subsubsection*{Lexicographic Unranking}

A companion technique to digit DP is \textbf{unranking}: given a count function, find the $n$th element in lexicographic order.

\begin{enumerate}
    \item For each prefix length $\ell$, count valid completions using the DP.
    \item Find the length $\ell^*$ such that the cumulative count first reaches or exceeds $n$.
    \item Build the number digit by digit: for each position, try each digit $d$ in order, count completions with that digit placed, subtract from $n$ until the correct branch is identified.
\end{enumerate}

\begin{example}[Very odd numbers]
Call a positive integer \emph{very odd} if it uses only odd digits $\{1,3,5,7,9\}$, is divisible by $105 = 3 \times 5 \times 7$, and each odd digit appears an odd number of times. Let $\Theta(n)$ be the $n$th such number.

\textbf{Key observations:}
\begin{itemize}
    \item Divisibility by 5 with only odd digits forces the last digit to be 5. Write the number as a prefix $P$ followed by 5.
    \item Divisibility by 7: $(10P + 5) \equiv 0 \pmod 7 \Rightarrow P \equiv 3 \pmod 7$.
    \item Divisibility by 3: total digit sum $\equiv 0 \pmod 3$. Since the appended 5 contributes 5, we need $\text{digitSum}(P) \equiv 1 \pmod 3$.
    \item Parity of digit counts: in the final number all five digits $\{1,3,5,7,9\}$ must appear an odd number of times. Appending a 5 flips the parity of 5, so in $P$: digits $1,3,7,9$ must appear an \emph{odd} number of times and digit 5 an \emph{even} number of times.
\end{itemize}

\textbf{DP state}: $(\text{rem},\; \text{mod}_7,\; \text{mod}_3,\; \text{parityMask})$ where $\text{parityMask}$ is a 5-bit integer tracking the parities of the five odd-digit counts. At $\text{rem}=0$ a state is valid iff $\text{mod}_7 = 3$, $\text{mod}_3 = 1$, and $\text{parityMask} = \texttt{11011}_2$ (bits for $1,3,7,9$ set; bit for 5 clear).

Results: $\Theta(1) = 1117935$, $\Theta(10^3) = 11137955115$, $\Theta(10^{16}) = 13313751171933973557517973175$.
\end{example}

\subsection{Binary Trie DP (Expected-Value Guessing)}

A \textbf{binary trie} stores a set of integers by inserting their binary representations as root-to-leaf paths. Reading bits from the \emph{least significant bit} (LSB) first gives a natural trie over the bit-decomposition of each number.

\subsubsection*{Optimal Guessing DP}

\begin{example}[Expected score when guessing a prime bit by bit]
A prime $p \le N$ is drawn uniformly at random. Its binary representation is revealed one bit at a time starting from the LSB; before each bit is shown the guesser chooses 0 or 1 to maximise expected correct guesses.

\textbf{Approach}: insert all primes $\le N$ into a binary trie (LSB-first). At each trie node $u$, let $c_0(u)$ and $c_1(u)$ be the number of candidates with next bit 0 or 1, and $t(u) = c_0 + c_1$. The optimal expected score from node $u$ is:
\[
V(u) = \frac{\max(c_0(u),\, c_1(u))}{t(u)}
+ \frac{t(\text{child}_0)}{t(u)}\,V(\text{child}_0)
+ \frac{t(\text{child}_1)}{t(u)}\,V(\text{child}_1).
\]
The first term is the probability of guessing the current bit correctly (always guess the majority). The remaining terms are the expected future score weighted by which branch is actually taken.

Compute $V$ bottom-up (leaves first). The answer is $V(\text{root})$.

\textbf{Sanity checks}: $E(10) = 2$, $E(30) = 2.9$, $E(10^8) = 14.97696693$.
\end{example}

\begin{remark}[Why majority guessing is optimal]
At any node the guesser sees only the count distribution, not which prime was drawn. Guessing the majority bit maximises the probability of a correct guess at that step. This greedy-per-step strategy is globally optimal because guessing correctly does not affect the subsequent state (the true bit is always revealed afterwards).
\end{remark}

% ==============================================================
% Graph Theory
% ==============================================================

\section{Graph Theory}

\subsection{Paths, cycles, connectivity}

% TODO

\subsection{Trees and spanning trees}

% TODO

\subsection{Planarity, Euler characteristic, graph colouring}

% TODO

% ==============================================================
% Number Theory (Discrete)
% ==============================================================

\section{Number Theory (Discrete)}

\subsection{GCD, Bezout, extended Euclidean}

% TODO

\subsection{Congruences and modular inverses}

% TODO

\subsection{RSA and applications to cryptography}

% TODO


% ================================================================
%  PART VI — Calculus  (MA1521 + MA2104)
% ================================================================
\part{Calculus}
% ==============================================================
% Limits and Continuity
% ==============================================================

\section{Limits and Continuity}

\subsection{Epsilon-delta definitions}

% TODO

\subsection{Limit laws, squeeze theorem}

% TODO

\subsection{Continuity and uniform continuity}

% TODO

% ==============================================================
% Single-Variable Differentiation
% ==============================================================

\section{Single-Variable Differentiation}

\subsection{Differentiation rules and implicit differentiation}

% TODO

\subsection{Mean value theorem, L'Hopital's rule}

% TODO

\subsection{Optimisation and curve sketching}

% TODO

% ==============================================================
% Single-Variable Integration
% ==============================================================

\section{Single-Variable Integration}

\subsection{Fundamental theorem of calculus}

% TODO

\subsection{Integration techniques: IBP, partial fractions, trig subs}

% TODO

\subsection{Improper integrals}

% TODO

% ==============================================================
% Sequences and Series (MA1521)
% ==============================================================

\section{Sequences and Series (MA1521)}

\subsection{Convergence tests: ratio, root, comparison}

% TODO

\subsection{Power series and radius of convergence}

% TODO

\subsection{Taylor and Maclaurin series}

% TODO

% ==============================================================
% Vectors and Curves in Space
% ==============================================================

\section{Vectors and Curves in Space}

\subsection{Parametric curves, arc length}

% TODO

\subsection{Curvature, TNB frame}

% TODO

\subsection{Vector-valued functions}

% TODO

% ==============================================================
% Partial Derivatives
% ==============================================================

\section{Partial Derivatives}

\subsection{Partial derivatives and gradients}

% TODO

\subsection{Directional derivatives}

% TODO

\subsection{Chain rule for multivariable functions}

% TODO

% ==============================================================
% Optimisation in Several Variables
% ==============================================================

\section{Optimisation in Several Variables}

\subsection{Critical points, second derivative test}

% TODO

\subsection{Constrained optimisation and Lagrange multipliers}

% TODO

\subsection{Bordered Hessians}

% TODO

% ==============================================================
% Multiple Integrals
% ==============================================================

\section{Multiple Integrals}

\subsection{Double and triple integrals}

% TODO

\subsection{Change of variables, Jacobians}

% TODO

\subsection{Polar, cylindrical, spherical coordinates}

% TODO

% ==============================================================
% Vector Calculus
% ==============================================================

\section{Vector Calculus}

\subsection{Line integrals and work}

% TODO

\subsection{Green's, Stokes', Divergence theorems}

% TODO

\subsection{Conservative fields and potentials}

% TODO


% ================================================================
%  PART VII — Linear Algebra  (MA2001 + MA2101)
% ================================================================
\part{Linear Algebra}
% ==============================================================
% Vector Spaces
% ==============================================================

\section{Vector Spaces}

\subsection{Axioms, subspaces, span}

% TODO

\subsection{Basis, dimension, coordinates}

% TODO

\subsection{Linear independence}

% TODO

% ==============================================================
% Linear Maps and Matrices
% ==============================================================

\section{Linear Maps and Matrices}

\subsection{Null space and column space}

% TODO

\subsection{Rank-nullity theorem}

% TODO

\subsection{Change of basis}

% TODO

% ==============================================================
% Systems and Gaussian Elimination
% ==============================================================

\section{Systems and Gaussian Elimination}

\subsection{Row reduction and REF/RREF}

% TODO

\subsection{Solution sets and parameterisation}

% TODO

\subsection{LU decomposition}

% TODO

% ==============================================================
% Determinants
% ==============================================================

\section{Determinants}

\subsection{Cofactor expansion, properties}

% TODO

\subsection{Adjugate and Cramer's rule}

% TODO

\subsection{Geometric interpretation}

% TODO

% ==============================================================
% Eigenvalues and Eigenvectors
% ==============================================================

\section{Eigenvalues and Eigenvectors}

\subsection{Characteristic polynomial}

% TODO

\subsection{Diagonalisation}

% TODO

\subsection{Jordan normal form (intro)}

% TODO

% ==============================================================
% Inner Product Spaces
% ==============================================================

\section{Inner Product Spaces}

\subsection{Inner products and norms}

% TODO

\subsection{Cauchy-Schwarz, triangle inequality}

% TODO

\subsection{Gram-Schmidt orthogonalisation}

% TODO

% ==============================================================
% Matrix Decompositions
% ==============================================================

\section{Matrix Decompositions}

\subsection{QR decomposition}

% TODO

\subsection{Singular Value Decomposition (SVD)}

% TODO

\subsection{Spectral theorem and positive definite matrices}

% TODO


% ================================================================
%  PART VIII — Mathematical Analysis  (MA2108)
% ================================================================
\part{Mathematical Analysis}
% ==============================================================
% The Real Number System
% ==============================================================

\section{The Real Number System}

\subsection{Ordered field axioms}

% TODO

\subsection{Completeness and least upper bound}

% TODO

\subsection{Archimedean property}

% TODO

% ==============================================================
% Sequences of Real Numbers
% ==============================================================

\section{Sequences of Real Numbers}

\subsection{Convergence, limit laws}

% TODO

\subsection{Monotone convergence, Bolzano-Weierstrass}

% TODO

\subsection{Cauchy sequences}

% TODO

% ==============================================================
% Limits and Continuity (Analysis)
% ==============================================================

\section{Limits and Continuity (Analysis)}

\subsection{Sequential and epsilon-delta definitions}

% TODO

\subsection{Uniform continuity}

% TODO

\subsection{Intermediate and extreme value theorems}

% TODO

% ==============================================================
% Differentiability
% ==============================================================

\section{Differentiability}

\subsection{Definition and basic theorems}

% TODO

\subsection{Mean value and Rolle's theorems}

% TODO

\subsection{Taylor's theorem with remainder}

% TODO

% ==============================================================
% Riemann Integration
% ==============================================================

\section{Riemann Integration}

\subsection{Darboux and Riemann sums}

% TODO

\subsection{Integrability criteria}

% TODO

\subsection{Fundamental theorem, change of variables}

% TODO

% ==============================================================
% Series of Real Numbers
% ==============================================================

\section{Series of Real Numbers}

\subsection{Convergence tests (rigorous)}

% TODO

\subsection{Absolute and conditional convergence}

% TODO

\subsection{Uniform convergence of function series}

% TODO


% ================================================================
%  PART IX — Probability and Statistics  (ST2334)
% ================================================================
\part{Probability and Statistics}
% ==============================================================
% Probability Axioms
% ==============================================================

\section{Probability Axioms}

\subsection{Sample spaces, events, sigma-algebras}

% TODO

\subsection{Axioms of probability}

% TODO

\subsection{Conditional probability, independence}

% TODO

% ==============================================================
% Discrete Random Variables
% ==============================================================

\section{Discrete Random Variables}

\subsection{PMF, CDF, expectation, variance}

% TODO

\subsection{Bernoulli, Binomial, Geometric, Poisson}

% TODO

\subsubsection*{Conditional Geometric Distribution}

A useful technique is computing the conditional distribution of a geometric-like stopping time given that one event occurs before another.

\begin{example}[First 6 before first 5]
Roll a fair 6-sided die repeatedly. Let $A$ = ``the first 6 appears before the first 5''. Find $E[N \mid A]$, where $N$ is the total number of rolls.

\textbf{Step 1 — Joint distribution.} On each roll independently, the outcome is: 6 with probability $1/6$, 5 with probability $1/6$, or ``other'' (1,2,3,4) with probability $4/6 = 2/3$. Event $A$ requires $N-1$ ``other'' rolls followed by a 6, so
\[
P(N = n \text{ and } A) = \left(\frac{4}{6}\right)^{n-1} \cdot \frac{1}{6}.
\]

\textbf{Step 2 — Find $P(A)$.}
\[
P(A) = \sum_{n=1}^{\infty} \left(\frac{4}{6}\right)^{n-1} \frac{1}{6}
= \frac{1/6}{1 - 4/6} = \frac{1/6}{2/6} = \frac{1}{2}.
\]

\textbf{Step 3 — Conditional PMF.}
\[
P(N = n \mid A) = \frac{P(N=n \text{ and } A)}{P(A)}
= \frac{(4/6)^{n-1}(1/6)}{1/2} = \left(\frac{2}{3}\right)^{n-1} \cdot \frac{1}{3}.
\]
This is a $\mathrm{Geometric}(1/3)$ distribution: each ``trial'' under the conditioning succeeds (6 appears) with probability $1/3$ and fails (other, not 5) with probability $2/3$.

\textbf{Step 4 — Conditional expectation.}
\[
E[N \mid A] = \frac{1}{1/3} = 3.
\]

\textbf{Intuition.} Conditioning on $A$ is equivalent to pretending the die has only two outcomes: roll a 6 (probability $1/3$ among $\{1,2,3,4,6\}$) or roll something else (probability $2/3$). The mean waiting time to the 6 in this reduced experiment is 3.
\end{example}

\begin{remark}[General pattern]
Suppose events $A$ (probability $p$) and $B$ (probability $q$) are mutually exclusive on each trial, with ``neither'' having probability $1-p-q$. Then $P(\text{A before B}) = p/(p+q)$, and given A occurs first, $N \mid (\text{A before B}) \sim \mathrm{Geometric}(p/(p+q))$ with mean $(p+q)/p$.
\end{remark}

\subsection{Generating functions (PGF, MGF)}

% TODO

% ==============================================================
% Continuous Random Variables
% ==============================================================

\section{Continuous Random Variables}

\subsection{PDF, CDF, quantiles}

% TODO

\subsection{Uniform, Exponential, Gamma, Normal}

% TODO

\subsection{Chi-square, $t$, and $F$ distributions}

Everything begins with the standard normal $Z \sim N(0,1)$.
All three distributions below are built from this single building block.

\subsubsection*{Chi-Square Distribution}

\begin{definition}[Chi-square distribution]
If $Z_1, Z_2, \ldots, Z_k \stackrel{\text{iid}}{\sim} N(0,1)$, then
\[
X = Z_1^2 + Z_2^2 + \cdots + Z_k^2 \sim \chi^2(k).
\]
The parameter $k$ is the \textbf{degrees of freedom} and equals the number of squared normals added.
\end{definition}

\begin{remark}[Why squaring matters]
A standard normal $Z$ is symmetric around 0 and can be negative. After squaring: (1) all values become $\geq 0$, and (2) large deviations in either direction collapse to large positive values (e.g.\ $(-3)^2 = 9$). This is why $\chi^2$ is always non-negative and right-skewed.
\end{remark}

\begin{remark}[Shape by degrees of freedom]
\begin{itemize}
    \item $k=1$: Spikes near zero then drops sharply — ``backwards J''. Most of the time $Z \approx 0$, so $Z^2 \approx 0$.
    \item $k=2,3,\ldots$: A hump develops. The mode is at $k-2$ (for $k \geq 2$); the mean is $k$.
    \item Large $k$: By the CLT, summing many independent squared normals approaches normality. $\chi^2(k)$ looks increasingly bell-shaped, but remains shifted right (never below 0).
\end{itemize}
\end{remark}

\begin{remark}[Additive property]
If $U \sim \chi^2(d_1)$ and $V \sim \chi^2(d_2)$ independently, then $U + V \sim \chi^2(d_1 + d_2)$. Degrees of freedom add.
\end{remark}

\begin{remark}[Sample variance connection]
For a sample of size $n$ from $N(\mu,\sigma^2)$:
\[
\frac{(n-1)S^2}{\sigma^2} \sim \chi^2(n-1).
\]
\end{remark}

\subsubsection*{$t$-Distribution}

\begin{definition}[$t$-distribution]
If $Z \sim N(0,1)$ and $V \sim \chi^2(\nu)$ are independent, then
\[
T = \frac{Z}{\sqrt{V/\nu}} \sim t(\nu).
\]
\end{definition}

\begin{remark}[Intuition]
The $t$-distribution arises when estimating a population mean without knowing $\sigma^2$. The sample variance $S^2$ is random, so the denominator $S/\sqrt{n}$ introduces extra uncertainty. The $\chi^2(\nu)$ in the denominator captures this randomness from variance estimation.
\end{remark}

\begin{remark}[Shape]
The $t$-distribution is symmetric (like the normal) but has heavier tails: occasionally the denominator $\sqrt{V/\nu}$ is small, making $T$ large. As $\nu \to \infty$, $V/\nu \to 1$ and $T \to N(0,1)$. In practice, $t \approx N(0,1)$ for large samples.
\end{remark}

\begin{theorem}[$t$ squared is $F$]
If $T \sim t(\nu)$, then $T^2 \sim F(1,\nu)$.
\end{theorem}

\subsubsection*{$F$-Distribution}

\begin{definition}[$F$-distribution]
If $U \sim \chi^2(d_1)$ and $V \sim \chi^2(d_2)$ are independent, then
\[
F = \frac{U/d_1}{V/d_2} \sim F(d_1, d_2).
\]
\end{definition}

\begin{remark}[Intuition]
The $F$-distribution compares two sources of variation. Dividing by degrees of freedom makes the pieces comparable regardless of how many squared normals went into each. It is used in ANOVA, comparing variances, and testing joint significance of regression coefficients.
\end{remark}

\begin{remark}[Shape]
$F > 0$ always (ratio of positive quantities). Right-skewed for the same reason as $\chi^2$. When $d_2 \to \infty$, the denominator stabilises to 1, so $F(d_1, \infty) \equiv \chi^2(d_1)/d_1$.
\end{remark}

\subsubsection*{The Full Family Tree}

\begin{center}
\begin{tabular}{lll}
$N(0,1) \xrightarrow{\text{square}}$ & $\chi^2(1) \xrightarrow{\text{add }k}$ & $\chi^2(k)$ \\[4pt]
$\dfrac{N(0,1)}{\sqrt{\chi^2(\nu)/\nu}}$ & $= t(\nu)$ & $\xrightarrow{\text{square}} F(1,\nu)$ \\[8pt]
$\dfrac{\chi^2(d_1)/d_1}{\chi^2(d_2)/d_2}$ & $= F(d_1, d_2)$ &
\end{tabular}
\end{center}

\begin{center}
\begin{tabular}{llll}
\textbf{Distribution} & \textbf{Built from} & \textbf{Shape} & \textbf{Used for} \\
\hline
$\chi^2(k)$ & Sum of $k$ squared normals & Right-skewed, $\geq 0$ & Variance, goodness of fit \\
$t(\nu)$ & Normal $\div$ scaled $\chi^2$ & Symmetric, heavy tails & Testing one coefficient \\
$F(d_1,d_2)$ & Ratio of two scaled $\chi^2$ & Right-skewed, $\geq 0$ & Testing multiple coefficients jointly \\
\end{tabular}
\end{center}

\begin{remark}[One-line intuitions]
\begin{itemize}
    \item $\chi^2$: \emph{how much total squared deviation is there?}
    \item $t$: \emph{is this one thing far from zero, accounting for noisy variance?}
    \item $F$: \emph{are these things jointly far from zero, as a ratio of variations?}
\end{itemize}
\end{remark}

\subsection{Transformations and change of variables}

% TODO

% ==============================================================
% Joint Distributions
% ==============================================================

\section{Joint Distributions}

\subsection{Joint, marginal, conditional distributions}

% TODO

\subsection{Covariance and correlation}

% TODO

\subsection{Bivariate normal distribution}

% TODO

% ==============================================================
% Limit Theorems
% ==============================================================

\section{Limit Theorems}

\subsection{Law of large numbers}

% TODO

\subsection{Central limit theorem}

% TODO

\subsection{Delta method}

% TODO

% ==============================================================
% Statistical Inference
% ==============================================================

\section{Statistical Inference}

\subsection{Point estimation: MLE, MOM}

\subsubsection*{Estimators and Estimates}

\begin{definition}
An \textbf{estimator} is a function of sample data used to infer an unknown population parameter $\theta$. Because it depends on random sample data, an estimator $\hat{\theta}$ is itself a \textbf{random variable} with its own probability distribution. The specific numerical value produced from a given sample is called an \textbf{estimate}.
\end{definition}

\begin{example}
The sample mean $\bar{X} = \frac{1}{n}\sum_{i=1}^n X_i$ is an \emph{estimator} of the population mean $\mu$. If a particular sample yields $\bar{x} = 172\,\text{cm}$, then $172\,\text{cm}$ is the \emph{estimate}.
\end{example}

\begin{remark}
Point estimators provide a single value; interval estimators (confidence intervals) provide a range of plausible values. See the next subsection for confidence intervals.
\end{remark}

\subsubsection*{Properties of Good Estimators}

\begin{definition}
Let $\hat{\theta}$ be an estimator of parameter $\theta$, based on a sample of size $n$.
\begin{itemize}
    \item \textbf{Unbiasedness:} $\hat{\theta}$ is \emph{unbiased} if $\mathbb{E}[\hat{\theta}] = \theta$. The \emph{bias} is $\mathrm{Bias}(\hat{\theta}) = \mathbb{E}[\hat{\theta}] - \theta$.
    \item \textbf{Consistency:} $\hat{\theta}$ is \emph{consistent} if $\hat{\theta} \xrightarrow{p} \theta$ as $n \to \infty$; i.e.\ the estimator converges in probability to the true value.
    \item \textbf{Efficiency:} Among all unbiased estimators, $\hat{\theta}$ is \emph{efficient} if it achieves the minimum variance. The \textbf{Cramér--Rao lower bound} provides the theoretical minimum variance for any unbiased estimator.
\end{itemize}
\end{definition}

\begin{example}
The sample mean $\bar{X}$ is unbiased for $\mu$. The sample variance $S^2 = \frac{1}{n-1}\sum_{i=1}^n (X_i - \bar{X})^2$ is unbiased for $\sigma^2$ (the $n-1$ denominator corrects for the bias that arises using $n$).
\end{example}

\subsubsection*{Method of Moments (MOM)}

\begin{definition}
The \textbf{Method of Moments} (MOM) estimates parameters by equating \emph{population moments} to \emph{sample moments} and solving. The $k$-th population moment is $\mu_k' = \mathbb{E}[X^k]$; the corresponding sample moment is $m_k' = \frac{1}{n}\sum_{i=1}^n x_i^k$.
\end{definition}

\begin{example}
For a distribution with mean $\mu$ and variance $\sigma^2$, MOM sets:
\[
\hat{\mu} = m_1' = \bar{x}, \qquad \hat{\sigma}^2 = m_2' - (m_1')^2 = \frac{1}{n}\sum x_i^2 - \bar{x}^2
\]
MOM estimators are simple to compute but may not always be the most efficient.
\end{example}

\subsubsection*{Maximum Likelihood Estimation (MLE)}

\begin{definition}
Given a sample $x_1, \ldots, x_n$ and a parametric model with parameter $\theta$, the \textbf{likelihood function} is:
\[
L(\theta) = \prod_{i=1}^n f(x_i \mid \theta)
\]
The \textbf{maximum likelihood estimator} (MLE) $\hat{\theta}_{\mathrm{MLE}}$ is the value of $\theta$ that maximises $L(\theta)$, or equivalently the \textbf{log-likelihood} $\ell(\theta) = \ln L(\theta)$ (since $\ln$ is monotone):
\[
\hat{\theta}_{\mathrm{MLE}} = \arg\max_{\theta}\; \ell(\theta)
\]
In practice, solve $\frac{\partial \ell}{\partial \theta} = 0$ (the \emph{score equation}).
\end{definition}

\begin{example}
For $X_1,\ldots,X_n \overset{\text{iid}}{\sim} \mathcal{N}(\mu, \sigma^2)$, the log-likelihood is:
\[
\ell(\mu, \sigma^2) = -\frac{n}{2}\ln(2\pi\sigma^2) - \frac{1}{2\sigma^2}\sum_{i=1}^n (x_i - \mu)^2
\]
Setting $\partial\ell/\partial\mu = 0$ gives $\hat{\mu}_{\mathrm{MLE}} = \bar{x}$. Setting $\partial\ell/\partial\sigma^2 = 0$ gives $\hat{\sigma}^2_{\mathrm{MLE}} = \frac{1}{n}\sum(x_i-\bar{x})^2$ (note: biased, unlike MOM's sample variance).
\end{example}

\begin{remark}
MLEs have strong asymptotic properties: they are consistent, asymptotically normal, and asymptotically efficient (achieving the Cramér--Rao bound as $n \to \infty$).
\end{remark}

\subsection{Confidence intervals}

% TODO

\subsection{Hypothesis testing: z, t, chi-squared, F tests}

\subsubsection*{The Chi-Squared ($\chi^2$) Test}

\begin{definition}
The \textbf{chi-squared test statistic} measures how far observed categorical frequencies deviate from expected frequencies:
\[
\chi^2 = \sum_{i} \frac{(O_i - E_i)^2}{E_i}
\]
where $O_i$ is the \emph{observed} count and $E_i$ is the \emph{expected} count in category $i$. Under the null hypothesis, this statistic follows a $\chi^2$ distribution with an appropriate number of degrees of freedom.
\end{definition}

\begin{remark}
A smaller $\chi^2$ value means observed data align closely with expectations (supporting $H_0$). A large $\chi^2$ value — falling in the critical region — provides evidence against $H_0$.
\end{remark}

\subsubsection*{Type I: Goodness-of-Fit Test}

\begin{definition}
A \textbf{goodness-of-fit test} asks: does a sample distribution match a hypothesised population distribution? With $k$ categories, the degrees of freedom are $\nu = k - 1$ (minus one additional for each estimated parameter).
\end{definition}

\begin{example}
A die is rolled 60 times. Under $H_0$ (fair die), each face is expected $E_i = 10$ times. The observed counts are $\{8, 12, 9, 11, 7, 13\}$. Then:
\[
\chi^2 = \frac{(8-10)^2}{10} + \frac{(12-10)^2}{10} + \cdots + \frac{(13-10)^2}{10} = \frac{4+4+1+1+9+9}{10} = 2.8
\]
With $\nu = 5$ degrees of freedom, compare to the critical value $\chi^2_{0.05,5} = 11.07$. Since $2.8 < 11.07$, we fail to reject $H_0$.
\end{example}

\subsubsection*{Type II: Test of Independence}

\begin{definition}
A \textbf{test of independence} determines whether two categorical variables are associated. Data are arranged in an $r \times c$ contingency table. The expected count for cell $(i,j)$ is:
\[
E_{ij} = \frac{(\text{row } i \text{ total}) \times (\text{column } j \text{ total})}{\text{grand total}}
\]
Degrees of freedom: $\nu = (r-1)(c-1)$.
\end{definition}

\begin{example}
To test whether gender and voting preference are independent, arrange counts into a $2 \times 3$ contingency table (2 genders, 3 parties). Compute $E_{ij}$ for each cell, then sum $\frac{(O_{ij}-E_{ij})^2}{E_{ij}}$ across all cells. With $\nu = (2-1)(3-1) = 2$, compare to $\chi^2_{0.05,2} = 5.99$.
\end{example}

\subsubsection*{Requirements and Limitations}

\begin{remark}
The $\chi^2$ test requires:
\begin{itemize}
    \item \textbf{Categorical data} — it does not apply to continuous measurements directly.
    \item \textbf{Independent observations} — each subject appears in exactly one cell.
    \item \textbf{Sufficient expected counts} — every cell should have $E_i \geq 5$. If this fails, merge categories or use Fisher's exact test.
\end{itemize}
The $\chi^2$ distribution is an approximation; it becomes accurate as sample size grows. For a $2 \times 2$ table with small samples, apply Yates' continuity correction:
\[
\chi^2 = \sum_{i} \frac{(|O_i - E_i| - 0.5)^2}{E_i}
\]
\end{remark}

\begin{remark}
For an application of the goodness-of-fit test to real-world fraud detection, see the Benford's Law subsection below, where the $\chi^2$ statistic is used to detect suspiciously uniform leading-digit distributions in financial data.
\end{remark}

\subsubsection*{The $F$-Test}

\begin{definition}[$F$-test]
An \textbf{$F$-test} is used to test joint hypotheses about multiple parameters simultaneously. Under the null, the test statistic follows an $F(q, n-k)$ distribution where $q$ is the number of restrictions and $n-k$ is the residual degrees of freedom.
\end{definition}

\begin{remark}[Why we need $F$ instead of repeated $t$-tests]
Testing $H_0\colon \beta_1 = 0$ \emph{and} $\beta_2 = 0$ with two separate $t$-tests overstates confidence: each test has a false-positive probability, and running two inflates the overall Type I error rate. The $F$-test tests both restrictions simultaneously and controls the error rate correctly.
\end{remark}

\begin{theorem}[Large-sample $F$-statistic for $q$ restrictions]
In large samples, when $q$ restrictions $H_0\colon R\beta = r$ are tested jointly, the $F$-statistic follows
\[
F \sim F(q, \infty).
\]
This is because each restricted $t$-statistic is approximately $N(0,1)$ under $H_0$, so their squares are $\chi^2(1)$ and their sum is $\chi^2(q)$:
\[
F = \frac{\chi^2(q)}{q} \sim F(q,\infty).
\]
\end{theorem}

\begin{remark}[Connection to the distributions in ch.\ 62]
$F(q,\infty)$ is the limit of $F(q, d_2)$ as $d_2 \to \infty$, at which point the denominator stabilises to 1 and $F = \chi^2(q)/q$. For regression in finite samples, the denominator uses the residual $\chi^2(n-k)$ in place of the infinite-df limit.
\end{remark}

\subsection{Benford's Law and Fraud Detection}

\begin{definition}
\textbf{Benford's Law} (also called the \textit{First-Digit Law}) states that in many naturally occurring datasets, the probability that the leading (first significant) digit is $d \in \{1, 2, \ldots, 9\}$ follows a logarithmic distribution:
\[
P(D = d) = \log_{10}\!\left(1 + \frac{1}{d}\right)
\]
\end{definition}

\begin{remark}
The predicted probabilities are approximately:
\begin{center}
\begin{tabular}{c|ccccccccc}
Digit $d$ & 1 & 2 & 3 & 4 & 5 & 6 & 7 & 8 & 9 \\
\hline
$P(D=d)$ & 30.1\% & 17.6\% & 12.5\% & 9.7\% & 7.9\% & 6.7\% & 5.8\% & 5.1\% & 4.6\%
\end{tabular}
\end{center}
Crucially, the digit \textbf{1} appears as the leading digit nearly a third of the time --- far more often than a naive uniform assumption ($11.\overline{1}\%$) would suggest.
\end{remark}

\subsubsection*{Logarithmic Interpretation}

The law arises from the \emph{gaps} between consecutive base-10 logarithms of integers:
\[
P(D = d) = \log_{10}(d+1) - \log_{10}(d) = \log_{10}\!\left(\frac{d+1}{d}\right) = \log_{10}\!\left(1+\frac{1}{d}\right)
\]
These gaps are largest at $d=1$ and shrink monotonically toward $d=9$, explaining why lower digits dominate.

\begin{theorem}[Condition for Benford's Law]
Benford's Law holds for a dataset when the \textbf{fractional part} of $\log_{10}(x)$ --- i.e.\ $\{\log_{10}(x)\}$ --- is uniformly distributed on $[0,1)$. Under this condition, the probability of each leading digit is exactly the logarithmic gap above.
\end{theorem}

\subsubsection*{Why Naturally Occurring Data Follows This Pattern}

Consider a quantity growing multiplicatively (e.g.\ compound interest at 10\% per year):
\begin{itemize}
    \item Going from \$1 to \$2 requires a $100\%$ increase (and takes many years).
    \item Going from \$8 to \$9 requires only an $11\%$ increase (reached quickly).
\end{itemize}
Numbers therefore spend \emph{proportionally more time} with leading digit 1 than with 9, regardless of the unit of measurement.

\begin{theorem}[Scale Invariance]
Benford's Law is \textbf{scale invariant}: multiplying an entire dataset by any positive constant (e.g.\ converting currencies or units) does not change the leading-digit distribution. This is a hallmark property that distinguishes naturally generated data from fabricated data.
\end{theorem}

\subsubsection*{Limitations}

Benford's Law \emph{fails} when:
\begin{itemize}
    \item The data spans a narrow range of magnitudes (e.g.\ adult human heights in metres, all clustered near $1$--$2$ m).
    \item The dataset is artificially constrained (e.g.\ prices set to always begin with \$9.99).
    \item Data are purely random (uniformly or normally distributed with no multiplicative structure).
\end{itemize}

\subsubsection*{Application: Fraud Detection}

\begin{example}
Given an unlabelled dataset of financial figures, one can test conformity to Benford's Law using a \textbf{chi-squared goodness-of-fit test}:
\[
\chi^2 = \sum_{d=1}^{9} \frac{(O_d - E_d)^2}{E_d}, \qquad E_d = N \cdot \log_{10}\!\left(1+\frac{1}{d}\right)
\]
where $O_d$ is the observed count of leading digit $d$ and $N$ is the total number of values. A statistically significant deviation (especially a distribution that is \emph{too uniform}) is a red flag for fabricated data, since humans inventing numbers tend to distribute digits more evenly than nature does.
\end{example}

\begin{remark}
Benford's Law has been used to audit elections, detect accounting fraud, and expose billion-dollar financial scandals. It is context-free: no knowledge of what the numbers represent is required, only that the data plausibly spans several orders of magnitude. Non-conformity is not definitive proof of fraud, but it is a statistically grounded signal warranting further investigation.
\end{remark}


% ================================================================
%  PART X — Stochastic Processes  (MA3238 + MA4251)
% ================================================================
\part{Stochastic Processes}
% ==============================================================
% Discrete-Time Markov Chains
% ==============================================================

\section{Discrete-Time Markov Chains}

\subsection{Transition matrices and Chapman-Kolmogorov}

% TODO

\subsection{Classification of states: recurrent, transient}

% TODO

\subsection{Stationary distributions and ergodicity}

% TODO

% ==============================================================
% Poisson Processes
% ==============================================================

\section{Poisson Processes}

\subsection{Definition and properties}

% TODO

\subsection{Superposition, thinning, colouring}

% TODO

\subsection{Compound Poisson process}

% TODO

% ==============================================================
% Continuous-Time Markov Chains
% ==============================================================

\section{Continuous-Time Markov Chains}

\subsection{Generator matrix and Kolmogorov equations}

% TODO

\subsection{Birth-death processes}

% TODO

\subsection{Embedded chains and limiting distributions}

% TODO

% ==============================================================
% Renewal Theory
% ==============================================================

\section{Renewal Theory}

\subsection{Renewal equation and elementary renewal theorem}

% TODO

\subsection{Alternating renewal processes}

% TODO

\subsection{Renewal reward theorem}

% TODO

% ==============================================================
% Martingales
% ==============================================================

\section{Martingales}

\subsection{Filtrations and conditional expectation}

% TODO

\subsection{Optional stopping theorem}

% TODO

\subsection{Martingale convergence}

% TODO

% ==============================================================
% Brownian Motion
% ==============================================================

\section{Brownian Motion}

\subsection{Definition and basic properties}

% TODO

\subsection{Reflection principle and first passage}

% TODO

\subsection{Quadratic variation}

% TODO

% ==============================================================
% Stochastic Calculus
% ==============================================================

\section{Stochastic Calculus}

\subsection{Ito integral and Ito's lemma}

% TODO

\subsection{Stochastic differential equations}

% TODO

\subsection{Girsanov's theorem (intro)}

% TODO


% ================================================================
%  PART XI — Mathematics for Reinforcement Learning  (MA4275)
% ================================================================
\part{Mathematics for Reinforcement Learning}
% ==============================================================
% Dynamic Programming
% ==============================================================

\section{Dynamic Programming}

\subsection{Principle of optimality}

% TODO

\subsection{Bellman equation (deterministic)}

% TODO

\subsection{Value iteration and policy iteration}

% TODO

% ==============================================================
% Markov Decision Processes
% ==============================================================

\section{Markov Decision Processes}

\subsection{State, action, reward, transition model}

% TODO

\subsection{Finite and infinite horizon problems}

% TODO

\subsection{Discounted and average reward criteria}

% TODO

% ==============================================================
% Value Functions and Policies
% ==============================================================

\section{Value Functions and Policies}

\subsection{State-value and action-value functions}

% TODO

\subsection{Optimal Bellman equations}

% TODO

\subsection{Fixed-point characterisation}

% TODO

% ==============================================================
% Optimisation for RL
% ==============================================================

\section{Optimisation for RL}

\subsection{Convex optimisation basics}

% TODO

\subsection{Gradient descent and projected gradient}

% TODO

\subsection{Stochastic approximation and SGD}

% TODO

% ==============================================================
% Function Approximation
% ==============================================================

\section{Function Approximation}

\subsection{Linear function approximation}

% TODO

\subsection{Projected Bellman equation}

% TODO

\subsection{Convergence theory and stability}

% TODO

% ==============================================================
% Theoretical RL Results
% ==============================================================

\section{Theoretical RL Results}

\subsection{PAC learning and sample complexity}

% TODO

\subsection{Regret bounds and online learning}

% TODO

\subsection{Policy gradient theorems}

% TODO


% ================================================================
%  PART XII — Quantum Computing  (CS4268)
% ================================================================
\part{Quantum Computing}
% ==============================================================
% Hilbert Spaces and Qubits
% ==============================================================

\section{Hilbert Spaces and Qubits}

\subsection{Qubits and Dirac notation}

% TODO

\subsection{Tensor products and multi-qubit systems}

\begin{definition}
The \textbf{tensor product} (written $\otimes$) combines two quantum systems into a joint system. If qubit 1 is in state $|a\rangle$ and qubit 2 is in state $|b\rangle$, their combined state is:
\[
|a\rangle \otimes |b\rangle
\]
This is usually abbreviated by juxtaposition: $|a\rangle|b\rangle = |ab\rangle$.
\end{definition}

\begin{example}[Two-qubit states]
Using the standard basis vectors $|0\rangle = \begin{bmatrix}1\\0\end{bmatrix}$ and $|1\rangle = \begin{bmatrix}0\\1\end{bmatrix}$, the tensor product is computed as the Kronecker product — multiply the scalar entries of the first vector by the entire second vector:
\[
|0\rangle \otimes |1\rangle = \begin{bmatrix}1\\0\end{bmatrix} \otimes \begin{bmatrix}0\\1\end{bmatrix} = \begin{bmatrix}1\cdot\begin{bmatrix}0\\1\end{bmatrix}\\[4pt] 0\cdot\begin{bmatrix}0\\1\end{bmatrix}\end{bmatrix} = \begin{bmatrix}0\\1\\0\\0\end{bmatrix}
\]
This 4-dimensional vector is the 2-qubit basis state $|01\rangle$. Similarly, $|1\rangle \otimes |0\rangle = |10\rangle = \begin{bmatrix}0\\0\\1\\0\end{bmatrix}$.

An $n$-qubit system lives in a $2^n$-dimensional Hilbert space — the state space grows exponentially with the number of qubits.
\end{example}

\begin{remark}[Tensor product for gates]
The $\otimes$ notation also applies to quantum gates. Writing $H^{\otimes 2} = H \otimes H$ means ``apply the Hadamard gate independently to each of two qubits''. More generally, $U^{\otimes n}$ applies gate $U$ in parallel to all $n$ qubits.
\end{remark}

\subsection{Inner products and measurement}

% TODO

\subsection{Circled operator notation}

\begin{remark}[Reference: circled operators]
Several ``circled'' binary operators appear in mathematics and quantum computing. Their meaning is context-dependent:

\begin{center}
\begin{tabular}{c|l|l}
Symbol & Name & Common meaning \\
\hline
$\oplus$ & Circled plus & XOR / addition mod 2 / direct sum \\
$\otimes$ & Circled times & Tensor product / Kronecker product \\
$\ominus$ & Circled minus & Formal difference / subtraction-like operation \\
$\oslash$ & Circled slash & Quotient-like / author-defined division operation \\
$\odot$ & Circled dot & Elementwise (Hadamard) product / composition \\
$\circ$ & Small circle & Function composition: $(f \circ g)(x) = f(g(x))$ \\
$\bigoplus$ & Big circled plus & Direct sum over many terms: $\bigoplus_{i=1}^n V_i$ \\
$\bigotimes$ & Big circled times & Tensor product over many terms: $\bigotimes_{i=1}^n |q_i\rangle$ \\
\end{tabular}
\end{center}

\textbf{Quick intuition:}
\begin{itemize}
    \item $\oplus$ — special addition (bit-flipping arithmetic)
    \item $\otimes$ — gluing systems together (combining spaces)
    \item $\ominus$ — special subtraction
    \item $\oslash$ — special division
    \item $\odot$ — special elementwise multiplication
    \item $\circ$ — function composition (not a circled symbol, but related in appearance)
\end{itemize}
Unlike the standard $+, -, \times$, circled symbols are \emph{not universal} — always check the author's definition when encountering them in a new context.
\end{remark}

\input{parts/p12_quantum/ch80_quantum_gates_circuits}
\input{parts/p12_quantum/ch81_quantum_algorithms_i}
\input{parts/p12_quantum/ch82_quantum_fourier_transform}
\input{parts/p12_quantum/ch83_shors_grovers}
\input{parts/p12_quantum/ch84_quantum_information_theory}


\end{document}
